% Chapter 2 entries SET .. XCOPY (AN-01).

% SET -- os/milton/command.c builtin_set, cmd_set_parse (list; NAME=VALUE,
% name upcased; NAME= clears, silent; NAME alone queries:
% "Environment variable not defined"); os/milton/env.c (512-byte arena:
% "Out of environment space"); command_repl (COMSPEC, PROMPT, PATH).
% Gates: test-env, test-command, test-interp-set.
\begin{DosEntry}{SET}{}{internal}
\Purpose Sets, removes or displays environment variables.
\Format \SyntaxLine{\kw{SET} \opt{\ph{name}\opt{\kw{=}\opt{\ph{string}}}}}
\Type{Internal}
\Remarks
The \emph{environment} is an area of memory that holds named strings, called
environment variables, for InitechDOS and for the programs you run.

Typed with no parameter, \cmd{SET} lists every variable, one to a line, as
\ph{name}\texttt{=}\ph{string}.

\Param{name=string} sets the variable \ph{name} to \ph{string}, replacing
any earlier value.  The name is changed to capital letters; the string is
kept exactly as you type it.

\Param{name=} with no string removes the variable.  Nothing is displayed.

\Param{name} alone displays that variable, or
\texttt{Environment variable not defined}.

InitechDOS sets three variables when it starts: \texttt{COMSPEC}, the
location of the command processor; \texttt{PROMPT}, the command prompt (see
\cmd{PROMPT}); and \texttt{PATH}, the directories searched for commands
(see Chapter~1, ``The Practice Diskette'').  A batch file can use the value
of a variable by writing its name between percent signs, as
\texttt{\%PATH\%} (Chapter~4).

The environment holds 512 characters.  If a new variable does not fit,
\cmd{SET} displays \texttt{Out of environment space} and the environment is
not changed.
\Example
To list the environment, add the variable \texttt{DEPT}, and list it
again:
\Shot{set}
To display one variable, remove it, and look for it again:
\Shot{setclr}
\end{DosEntry}

% initech-2kgh: sort_program.asm, existing m0dc ascending/reverse goldens.
% Ref: MS-DOS 3.3 User's Reference p.116/PDF201.
\Gates{test-sort-filter, test-sort-filter-mutant, test-dos-folder}
\begin{DosEntry}{SORT}{}{external}
\Purpose Sorts input lines and displays them in order.
\Format \SyntaxLine{\kw{SORT} \opt{\kw{/R}} \kw{<} \ph{filename}}
\Type{External}
\Remarks
Input may come from a file or a pipe. Letters are sorted without regard
to capital or small letters. \Param{/R} reverses the order. Use \type{>}
and a file name to save the sorted output.
\Example
To save sorted lines, type \type{SORT < PHONE.TXT > SORTED.TXT}.
\end{DosEntry}

% SYS -- os/milton/diskutil.c sys_run (MS-DOS 3.3 UG p.119: the two system
% files, in order, into the first two root entries, fragmented sources
% allowed, COMMAND.COM NOT copied; the destination BPB kept, the InitechDOS
% boot code written; any failure "No room for system on destination
% disk"; success "System transferred").  Source = drive C (DEC-18).
% WL-0147.  Shipped as DOS\SYS.COM.  Gates: test-formatsys,
% test-formatsys-emu, test-format-system-boot.
\begin{DosEntry}{SYS}{System}{external}
\Purpose Transfers the InitechDOS system files to a disk.
\Format \SyntaxLine{\opt{\ph{d}\kw{:}}\opt{\ph{path}}\kw{SYS} \ph{d}\kw{:}}
\Type{External}
\Remarks
\Param{d:} is the drive of the disk to receive the system; it is required.
\cmd{SYS} copies the two system files, \file{IO.SYS} and \file{INITDOS.SYS},
from drive~C to the first two entries of the disk's root directory, and
writes the InitechDOS start-up code into its boot record.  The files already
on the disk are kept.  \cmd{SYS} does not copy \file{COMMAND.COM}; copy it
with \cmd{COPY}.

The first two root directory entries must be free, or already hold the
system files, as on a disk formatted with \kw{FORMAT /B}.  When the system
has been transferred, \cmd{SYS} displays \texttt{System transferred}.  If
the entries are in use, the disk has not enough free space, or the system
files cannot be read, it displays \texttt{No room for system on destination
disk} and the disk is not changed.
\Example
To transfer the system to the disk in drive~D, type \type{SYS D:}.
\end{DosEntry}

% TIME -- os/milton/command.c builtin_time, cmd_parse_time (hh:mm[:ss],
% 0..23 / 0..59; "Invalid time" then the unchanged time),
% cmd_format_time ("Current time is hh:mm:ss.cc").  Gates: test-datetime,
% test-command.
\begin{DosEntry}{TIME}{}{internal}
\Purpose Sets or displays the system time.
\Format \SyntaxLine{\kw{TIME} \opt{\ph{hh}\kw{:}\ph{mm}\opt{\kw{:}\ph{ss}}}}
\Type{Internal}
\Remarks
\Param{hh} is the hour, 0 through 23, on the 24-hour clock.

\Param{mm} is the minutes, 0 through 59.

\Param{ss} is the seconds, 0 through 59.  If you leave them out, the
seconds are set to 0.

With a valid time, \cmd{TIME} sets the system clock and displays the new
time.  With no parameter it displays the current time:
\texttt{Current time is} \ph{hh}\texttt{:}\ph{mm}\texttt{:}\ph{ss}%
\texttt{.}\ph{cc}, where \ph{cc} is hundredths of a second.

If the time is not valid, \cmd{TIME} displays \texttt{Invalid time}, leaves
the clock unchanged and displays the current time.
\Example
To set the clock to a quarter to two in the afternoon:
\Shot{time}
\end{DosEntry}

% TREE -- os/milton/xcopy_tree.c du_main(tree) (one path operand, default
% the root; /F; volume heading; "Path:", "Sub-directories:", ":" at column
% 16, "None", "Files:"; messages).  Shipped as APPS\TREE.COM.
% Gates: test-xcopy-tree, test-xcopy-tree-emu.
\begin{DosEntry}{TREE}{}{external}
\Purpose Displays the directories of a disk and, optionally, the files in
each.
\Format \SyntaxLine{\opt{\ph{d}\kw{:}}\opt{\ph{path}}\kw{TREE}
  \opt{\ph{d}\kw{:}}\opt{\ph{path}} \opt{\kw{/F}}}
\Type{External}
\Remarks
\Param{path} names the directory to start from.  If you leave it out,
\cmd{TREE} starts at the root directory.

\Param{/F} also lists the files in each directory.

\cmd{TREE} displays the heading \texttt{DIRECTORY PATH LISTING FOR VOLUME}
and the volume label.  Then, for the starting directory and each directory
below it, it displays \texttt{Path:} and the path, and
\texttt{Sub-directories:} followed by the names of the subdirectories, one
to a line, or \texttt{None}.  With \kw{/F}, \texttt{Files:} and the names
of the files follow, or \texttt{None}.

If the path is not a directory, \cmd{TREE} displays \texttt{Invalid path
(or file not found)}.  A parameter other than \kw{/F} displays
\texttt{Invalid parameter(s)}, and a second path displays \texttt{Invalid
number of parameters}.
\Example
To display the directories of the practice diskette:
\Shot{tree}
To list the files in the directory \file{SAMPLES}:
\Shot{treef}
\end{DosEntry}

% TYPE -- os/milton/command.c builtin_type (first token; open by name, no
% wildcard search; MSG-DOS-0003, MSG-DOS-0011; CR LF after the file).
% Gates: test-command, test-shell.
\begin{DosEntry}{TYPE}{}{internal}
\Purpose Displays the contents of a file.
\Format \SyntaxLine{\kw{TYPE} \opt{\ph{d}\kw{:}}\opt{\ph{path}}\ph{filename}\opt{\kw{.}\ph{ext}}}
\Type{Internal}
\Remarks
\Param{filename} names the file to display.  Wildcards are not allowed.

The file is displayed exactly as it is stored.  Use \cmd{TYPE} for text
files only; a program file contains codes that are not meant to be
displayed.

If the file does not exist, \cmd{TYPE} displays \texttt{File not found}.
With no file name it displays \texttt{Required parameter missing}.
\Example
To display the telephone list:
\Shot{type}
\end{DosEntry}

% VER -- os/milton/command.c builtin_ver (AH=30h: 3.30; "InitechDOS
% Version 3.30" after a blank line); spec/dos_banner.txt.  Gates:
% test-command, test-shell.
\begin{DosEntry}{VER}{Version}{internal}
\Purpose Displays the version number of InitechDOS.
\Format \SyntaxLine{\kw{VER}}
\Type{Internal}
\Remarks
\cmd{VER} displays a blank line and then the name and version of the
operating system.  Parameters are ignored.
\Example
\Shot{ver}
\end{DosEntry}

% VOL -- os/milton/command.c builtin_vol (drive A only, MSG-DOS-0004;
% extra parameters MSG-DOS-0010; MSG-DOS-0021/0022; read failure
% MSG-DOS-0009).  Gates: test-label-vol, test-dir-volume.
\begin{DosEntry}{VOL}{Volume}{internal}
\Purpose Displays the volume label of a disk.
\Format \SyntaxLine{\kw{VOL} \opt{\ph{d}\kw{:}}}
\Type{Internal}
\Remarks
\cmd{VOL} displays \texttt{Volume in drive A is} and the label, or
\texttt{Volume in drive A has no label}.  Create or change the label with
\cmd{LABEL}.

A drive other than A displays \texttt{Invalid drive specification}.
Anything after the drive displays \texttt{Too many parameters}.  If the disk
cannot be read, \cmd{VOL} displays \texttt{Access denied}.
\Example
\Shot{vol}
\end{DosEntry}

% XCOPY -- os/milton/xcopy_tree.c du_main (operands; /A /D:date /E /M /P /S
% /V /W; F/D question; names copied; "n File(s) copied"; exit codes 0, 1,
% 4, 5; cyclic and self-copy refusals; date/time kept via AH=57h; copies
% get archive set, source read-only kept).  Relative source from a
% subdirectory resolves against the root in the shipped build (no bead
% yet; WL-0119).  Shipped as APPS\XCOPY.COM.  Gates: test-xcopy-tree,
% test-xcopy-tree-emu, test-xcopy-tree-emu-cancel.
\begin{DosEntry}{XCOPY}{Extended Copy}{external}
\Purpose Copies groups of files, and the directories that contain them.
\Format \SyntaxLine{\opt{\ph{d}\kw{:}}\opt{\ph{path}}\kw{XCOPY}
  \opt{\ph{d}\kw{:}}\opt{\ph{path}}\opt{\ph{filename}\opt{\kw{.}\ph{ext}}}
  \opt{\ph{d}\kw{:}}\opt{\ph{path}}\opt{\ph{filename}\opt{\kw{.}\ph{ext}}}
  \opt{\kw{/A}\texttt{|}\kw{/M}} \opt{\kw{/D:}\ph{date}} \opt{\kw{/E}}
  \opt{\kw{/P}} \opt{\kw{/S}} \opt{\kw{/V}} \opt{\kw{/W}}}
\Type{External}
\Remarks
The first file specification names the source files.  It may be a file
name, a name with wildcards, or a directory, which means every file in it.

The second names the target.  If you leave it out, the files are copied to
the current directory.  If the target ends with a backslash, or is an
existing directory, the files are copied into it, and the directory is
created if necessary.  If the target does not exist and does not end with a
backslash, \cmd{XCOPY} asks:
\begin{TypeIn}
Does name specify a file name or directory name on the target
(F = file, D = directory)?
\end{TypeIn}
Press \type{F} to copy into a file of that name, or \type{D} to create a
directory.

\cmd{XCOPY} displays the name of each file as it copies it, and at the end
the number of files copied.  Each copy keeps the date and time of its
source and has its archive attribute set; a read-only source gives a
read-only copy.

\Param{/A} copies only files whose archive attribute is set, and leaves the
attribute unchanged.

\Param{/M} copies only files whose archive attribute is set, and clears the
attribute of each source file copied.  Use \kw{/M} to copy only the files
that have changed since the last copy.

\Param{/D:date} copies only files changed on or after \ph{date}, given as
for the \cmd{DATE} command.

\Param{/S} also copies the files in every subdirectory of the source,
creating the matching directories in the target.  Directories that hold no
files to copy are not created.

\Param{/E} used with \kw{/S} also creates the empty subdirectories.
\kw{/E} without \kw{/S} displays \texttt{Invalid parameter(s)}.

\Param{/P} asks before each file:
\ph{filename} \texttt{(Y/N?)}.  Press \type{Y} to copy the file or \type{N}
to skip it.

\Param{/V} reads back each copy and compares it with the source.

\Param{/W} displays \texttt{Press any key when ready to start copying
files} and waits, so that you can change diskettes first.

\cmd{XCOPY} will not copy a directory into one of its own subdirectories:
it displays \texttt{Cannot perform a cyclic copy}.  A file copied onto
itself displays \texttt{File cannot be copied onto itself}.

\cmd{XCOPY} returns exit code 0 when it copied files, 1 when no files were
found to copy, 4 for an error in the command line or a path, and 5 when a
file could not be read or written.

\draft{see bead initech-ktqk: run from a subdirectory, the shipped XCOPY finds a
relative source name in the root directory (XCOPY *.BAT in A:\bs SAMPLES
copies A:\bs AUTOEXEC.BAT); DOS 3.3 uses the current directory}
A source name without a path refers to the current directory.
\Example
The first command copies the text files into a new directory,
\file{BACKUP}; the second copies the directory \file{SAMPLES} to a new
directory, \file{SAVED}, the answer \type{D} being typed at the question:
\Shot{xcopy}
To be asked about each file, use \kw{/P}; here the answers are \type{Y},
\type{N} and \type{Y}:
\Shot{xcopyp}
\end{DosEntry}
